\documentclass[11pt,a4paper]{article}

\usepackage[utf8]{inputenc}
\usepackage[T1]{fontenc}
\usepackage{lmodern}
\usepackage[a4paper,margin=2.4cm]{geometry}
\usepackage{amsmath,amssymb,amsfonts}
\usepackage{bm}
\usepackage{graphicx}
\usepackage{hyperref}
\usepackage{authblk}
\usepackage{cite}
\usepackage{physics}
\usepackage{mathtools}
\usepackage{tensor}
\usepackage{enumitem}

\title{\textbf{RDCC--Euclid Interface (v49.1)}\\[0.2cm]
\large A Survey-Ready Framework for Testing Relaxation-Driven Cosmology}

\author[1]{Michael Lehmann}
\affil[1]{\small Independent Researcher, Relaxation-Driven Cyclic Cosmology (RDCC)\\
\texttt{mi.lehmann@gmx.de}}

\date{September 2026}

\begin{document}
\maketitle

% ========================================================
% ABSTRACT
% ========================================================

\begin{abstract}
Euclid provides high-precision weak-lensing and galaxy-clustering measurements
across a wide range of angular scales and redshifts. Relaxation-Driven Cyclic
Cosmology (RDCC) predicts geometric, morphological, and topological signatures
in these observables, all governed by the single infrared relaxation parameter
\(\alpha_{\mathrm{IR}}\). In the Phase-1 integrated framework these signatures
share one geometric origin: the freeze-out scale
\(k_{\mathrm{rel}}=a_{\mathrm{fr}}H(a_{\mathrm{fr}})\) derived in Companion~XII.
This Companion provides a survey-ready interface between RDCC and Euclid:
(i)~a modernized v49.1 theoretical overview, (ii)~the mapping from RDCC
signatures to Euclid observables, including growth, lensing, gravitational
slip and topology, (iii)~Euclid-level forecasts for \(\alpha_{\mathrm{IR}}\),
(iv)~degeneracy structure and multi-probe synergies, and (v)~a practical
integration pathway for Euclid pipelines. A detailed Nomenclature section is
included to ensure readability for Euclid scientists.
\end{abstract}

% ========================================================
% NOMENCLATURE
% ========================================================

\section*{Nomenclature}

\begin{tabular}{ll}
\(\alpha_{\mathrm{IR}}\) & Infrared relaxation parameter; single cosmological
parameter of RDCC \\[4pt]
\(a_{\mathrm{fr}}\) & Freeze-out scale factor of the relaxation sector
(\(\simeq 10^{-4}\)) \\[4pt]
\(k_{\mathrm{rel}}=a_{\mathrm{fr}}H(a_{\mathrm{fr}})\) &
Characteristic freeze-out scale governing growth suppression,
tensor damping, and gravitational slip \\[4pt]
\(\Phi\) & Gravitational potential of the visible sector \\[4pt]
\(\Psi\) & Newtonian-gauge potential sourcing peculiar velocities \\[4pt]
\(\Phi_{-}\) & Diffuse gravitational potential of the mirror sector \\[4pt]
\(\Phi_{\mathrm{tot}}\) & Total lensing potential including diffuse shadow \\[4pt]
\(\eta=\Phi/\Psi\) & Gravitational slip \\[4pt]
\(\kappa\) & Weak-lensing convergence field \\[4pt]
\(\gamma\) & Shear field \\[4pt]
\(P_{\kappa}(\ell)\) & Convergence power spectrum \\[4pt]
\(\mathrm{MF}_{i}\) & Minkowski functionals (area, boundary length,
Euler characteristic) \\[4pt]
\(\mathcal{D}\) & Persistence diagram (persistent homology) \\[4pt]
\(W_{2}\) & Optimal-transport distance between persistence diagrams \\[4pt]
\(f\sigma_8\) & Growth-rate observable \\[4pt]
\(\Delta N_{\mathrm{eff}}\) & Dark-radiation excess \\[4pt]
\(A_L\) & CMB-lensing amplitude \\[4pt]
\(\sigma(\alpha_{\mathrm{IR}})\) & Forecasted uncertainty on
\(\alpha_{\mathrm{IR}}\) \\[4pt]
\(\Delta\log Z\) & Bayesian evidence contrast \\[4pt]
\end{tabular}

% ========================================================
% 1 INTRODUCTION
% ========================================================

\section{Introduction}

Euclid's combination of weak lensing, galaxy clustering, tomography, and
topological probes makes it uniquely suited to test RDCC. The theory predicts
a set of correlated signatures across all Euclid observables, controlled by the
single parameter \(\alpha_{\mathrm{IR}}\). These signatures arise from:
\begin{itemize}
  \item the diffuse gravitational shadow of the mirror sector,
  \item scale-dependent suppression of the growth rate,
  \item a scale-dependent gravitational slip,
  \item mild modifications of the lensing potential,
  \item geometric and topological distortions in the convergence field.
\end{itemize}
They are not independent phenomenological effects. They share the freeze-out
scale of Companion~XII. This Companion provides a modernized v49.1 interface
between RDCC and Euclid.

% ========================================================
% 2 RDCC OVERVIEW (v49.1)
% ========================================================

\section{RDCC Overview (v49.1)}

RDCC describes the Universe as a globally coherent CPT-symmetric quantum state
with two thermodynamic sectors. The infrared relaxation parameter
\(\alpha_{\mathrm{IR}}\) quantifies the residual mismatch between the reduced
density matrices of the sectors after the bounce.

RDCC predicts:
\begin{itemize}
  \item a diffuse, non-clustering gravitational potential \(\Phi_{-}\),
  \item scale-dependent suppression of \(f\sigma_8\) above \(k_{\mathrm{rel}}\),
  \item a gravitational slip \(\eta=\Phi/\Psi\) that becomes scale-dependent
        near \(k_{\mathrm{rel}}\),
  \item mild enhancement of the lensing amplitude \(A_L\),
  \item geometric and topological distortions in \(\kappa\).
\end{itemize}
These signatures are directly accessible to Euclid.

% ========================================================
% 3 GEOMETRIC ORIGIN OF EUCLID SIGNATURES
% ========================================================

\section{Geometric Origin of Euclid Signatures}

The observable signatures of RDCC are not independent phenomenological
effects. They originate from a common geometric sequence,
\begin{equation}
  \alpha
  \;\rightarrow\;
  a_{\mathrm{fr}}
  \;\rightarrow\;
  k_{\mathrm{rel}}
  \;\rightarrow\;
  \{D,\, f,\, \gamma,\, \eta,\, \Omega_{\mathrm{GW}}\}.
\end{equation}
The relaxation scale
\begin{equation}
  k_{\mathrm{rel}}
  =
  a_{\mathrm{fr}} H(a_{\mathrm{fr}})
\end{equation}
was derived in Companion~XII from the freeze-out geometry of the relaxation
sector. Euclid therefore probes the same geometric transition that appears in
growth suppression (Companions~VII and~L), tensor damping (Companion~XV), and
gravitational slip (Companion~XLIV).

% ========================================================
% 4 RDCC SIGNATURES IN EUCLID OBSERVABLES
% ========================================================

\section{RDCC Signatures in Euclid Observables}

\subsection{Lensing Potential}

The total lensing potential is
\begin{equation}
  \Phi_{\mathrm{tot}}
  =
  \Phi + \Phi_{-},
\end{equation}
where \(\Phi_{-}\) is smooth and diffuse. In the RDCC framework the
mirror-sector potential \(\Phi_{-}\) obeys a Poisson equation of the form
\begin{equation}
  \nabla^{2} \Phi_{-}
  \propto
  4\pi G\, \rho_{(-)}\, \delta_{(-)},
\end{equation}
as developed in Companion~VII, with the diffuse, non-clustering character of
\(\rho_{(-)}\delta_{(-)}\) leading to a smooth contribution to the total
lensing potential. This produces a mild enhancement in \(P_{\kappa}(\ell)\)
and shifts peak and void statistics.

\subsection{Growth-Rate Suppression}

RDCC predicts
\begin{equation}
  f\sigma_8(k)
  =
  f\sigma_8^{\Lambda\mathrm{CDM}}
  \left[
    1 - \alpha_{\mathrm{IR}}\, \Theta(k - k_{\mathrm{rel}})
  \right],
\end{equation}
where the Heaviside step function \(\Theta\) represents an idealized
transition around the relaxation scale \(k_{\mathrm{rel}}\).

The transition scale is fixed by
\begin{equation}
  k_{\mathrm{rel}}
  =
  a_{\mathrm{fr}} H(a_{\mathrm{fr}})
\end{equation}
and is therefore inherited from the freeze-out geometry rather than
introduced phenomenologically.

For Euclid-level implementations a smoothed transition function can be used,
for example
\begin{equation}
  f\sigma_8(k)
  =
  f\sigma_8^{\Lambda\mathrm{CDM}}
  \left[
    1
    -
    \alpha_{\mathrm{IR}}\,
    \frac{1}{2}
    \left(
      1 + \tanh\frac{k - k_{\mathrm{rel}}}{\Delta k}
    \right)
  \right],
\end{equation}
with a characteristic width \(\Delta k\) controlling the smoothness of the
transition. Euclid's spectroscopic sample is highly sensitive to this
signature.

\subsection{Gravitational Slip}

Companion~XLIV established that the gravitational slip arises from an
effective anisotropic stress generated by the mirror sector and the
relaxation geometry. The slip parameter
\begin{equation}
  \eta
  =
  \frac{\Phi}{\Psi}
\end{equation}
becomes scale dependent near
\begin{equation}
  k_{\mathrm{rel}}.
\end{equation}
Consequently, weak-lensing observables (\(\Phi+\Psi\)) and clustering or
redshift-space-distortion observables (\(\Psi\)) provide complementary
probes of the same freeze-out transition.

\subsection{Geometric and Morphological Distortions}

RDCC induces characteristic distortions in:
\begin{itemize}
  \item peak counts,
  \item void distributions,
  \item Minkowski functionals,
  \item persistent homology diagrams.
\end{itemize}
These probes are robust against shape noise and survey systematics.

\subsection{Topological Rearrangements}

Persistent homology and optimal-transport distances \(W_{2}\) capture global
topological rearrangements induced by \(\Phi_{-}\).

% ========================================================
% 5 EUCLID-LEVEL FORECASTS
% ========================================================

\section{Euclid-Level Forecasts}

Euclid forecasts yield
\begin{equation}
  \sigma(\alpha_{\mathrm{IR}})
  \sim
  \mathcal{O}(10^{-2}),
\end{equation}
with strongest sensitivity from:
\begin{itemize}
  \item weak-lensing tomography,
  \item peak and void statistics,
  \item persistent homology,
  \item transport-based topology,
  \item combined multi-probe likelihoods.
\end{itemize}
The diffuse shadow contributes significantly to the Fisher information. The
same forecasts constrain the location of the freeze-out scale
\(k_{\mathrm{rel}}=a_{\mathrm{fr}}H(a_{\mathrm{fr}})\) once a numerical
evaluation of that product on the Companion-XI background is available.

% ========================================================
% 6 DEGENERACY STRUCTURE AND SYNERGY
% ========================================================

\section{Degeneracy Structure and Synergy}

RDCC variations can partially mimic shifts in \(\Lambda\)CDM parameters.
Euclid breaks these degeneracies through:
\begin{itemize}
  \item tomographic binning,
  \item cross-correlation with galaxy clustering,
  \item topological probes orthogonal to \(\Lambda\)CDM manifolds,
  \item multi-probe likelihoods.
\end{itemize}

The same relaxation scale is simultaneously constrained by
\begin{itemize}
  \item growth suppression,
  \item weak lensing,
  \item gravitational slip,
  \item topological observables,
  \item tensor observables.
\end{itemize}
Growth and RSD respond primarily to \(\Psi\); lensing responds to
\((\Phi+\Psi)/2\); topology responds to the diffuse shadow \(\Phi_{-}\);
the SGWB of Companion~XV responds to the same \(k_{\mathrm{rel}}\) through
tensor damping. This multi-probe structure is the operational content of the
geometric chain above.

% ========================================================
% 7 INTEGRATION INTO EUCLID PIPELINES
% ========================================================

\section{Integration into Euclid Pipelines}

RDCC can be integrated into Euclid's analysis framework through:
\begin{itemize}
  \item \textbf{Forward modeling:}
        modified lensing potential \(\Phi_{\mathrm{tot}}\) and a
        scale-dependent slip \(\eta(k)\) in simulations.
  \item \textbf{Summary statistics:}
        extension of Euclid's pipeline to include peaks, Minkowski
        functionals, persistent homology, and transport topology.
  \item \textbf{Likelihood integration:}
        modular addition of RDCC likelihood components sharing one
        \(k_{\mathrm{rel}}\) and one \(\alpha_{\mathrm{IR}}\).
  \item \textbf{Validation:}
        RDCC Benchmark Suite for simulation-based calibration.
\end{itemize}

% ========================================================
% 8 CONCLUSION
% ========================================================

\section{Conclusion}

RDCC provides a theoretically motivated, observationally testable extension
of \(\Lambda\)CDM. Its Euclid signatures are not a list of independent
effects. They descend from
\begin{equation}
  \alpha
  \;\rightarrow\;
  a_{\mathrm{fr}}
  \;\rightarrow\;
  k_{\mathrm{rel}}
  \;\rightarrow\;
  \{D,\, f,\, \gamma,\, \eta,\, \Omega_{\mathrm{GW}}\}.
\end{equation}
Euclid-level forecasts demonstrate percent-level sensitivity to
\(\alpha_{\mathrm{IR}}\) and can, once \(a_{\mathrm{fr}}H(a_{\mathrm{fr}})\)
is evaluated numerically, test whether the observed transition scale
coincides with the freeze-out geometry of Companion~XII. This Companion
provides a survey-ready interface for Euclid collaborators and a practical
pathway for testing relaxation-driven cosmology.

\end{document}